% =============================================================================
% cap5 §5.1 Posicionamento e leitura dos resultados
% =============================================================================
\section{Posicionamento}
\label{sec:resultados-posicionamento}

\obsfic{Os números, tabelas e gráficos deste capítulo são estimativas de faixa, extraídas da literatura recente dos modelos e datasets adotados, e servem de \emph{template} para a substituição direta pelos valores medidos após a execução do protocolo do Capítulo~\ref{cap:metodologia}.}

Os resultados aqui apresentados foram derivados de três fontes: as faixas de desempenho reportadas pelos próprios mantenedores dos modelos pré-treinados~\cite{shinoda2024petface,adam2024wildlifereid10k,akbar2025bodypart}, a literatura recente de classificação multi-espécie em \emph{camera trapping}~\cite{tabak2019nacti,norouzzadeh2018automatically} e a curva empírica de eficiência reportada por \citeauthor{velez2022choosing} (\citeyear{velez2022choosing}) para pipelines de duas etapas. As faixas mantêm coerência com o regime de avaliação em distribuição próxima da treinada (dados do portal LILA~BC) e devem ser lidas como \textbf{limite superior} do desempenho operacional esperado no Campus~2, conforme discutido na Seção~\ref{sec:metodologia-datasets}.

O capítulo organiza-se na ordem das fases (Seções~\ref{sec:resultados-fase2}--\ref{sec:resultados-fase4}), seguida do teste integrador (Seção~\ref{sec:resultados-end2end}) e da discussão consolidada (Seção~\ref{sec:resultados-discussao}). Em todos os casos, o reporte combina uma tabela quantitativa (faixas estimadas) e, quando aplicável, uma figura conceitual. Os comparativos com os baselines seguem a classificação qualitativa fixada na Seção~\ref{sec:metodologia-metricas}: diferença pequena, moderada ou grande.
